% === Begin Label Data ===
% ID: 566
% 难度星级: 
% 标签: 
% 备注: 
% 组卷引用次数: 0
% === End  Label Data ===

\begin{problem}{2013}{G}{湖南卷}{20}{解析几何}
在平面直角坐标系 $ xOy $ 中，将从点 $ M $ 出发沿纵、横方向到达点 $ N $ 的任一路径称为 $ M $ 到 $ N $ 的一条“$ L $ 路径”. 如图所示的路径 $ MM_1M_2M_3N $ 与路径 $ MN_1N $ 都是 $ M $ 到 $ N $ 的“$ L $ 路径”. 某地有三个新建的居民区，分别位于平面 $ xOy $ 内三点 $ A(3,20) $, $ B(-10,0) $, $ C(14,0) $ 处. 现计划在 $ x $ 轴上方区域（包含 $ x $ 轴）内的某一点 $ P $ 处修建一个文化中心.

\begin{tikzpicture}[line cap=round,line join=round,scale=0.75]
% axes
\draw[->] (-0.8,0) -- (7.6,0) node[below] {$x$};
\draw[->] (0,-0.5) -- (0,5.2) node[left] {$y$};
\node[below left] at (0,0) {$O$};

% polyline path: M -> M1 -> M2 -> M3 -> N -> N1 -> back to M (as in the sketch)
\coordinate (M)  at (1.6,3.6);
\coordinate (M1) at (5.6,3.6);
\coordinate (M2) at (5.6,2.4);
\coordinate (M3) at (4.4,2.4);
\coordinate (N)  at (4.4,1.6);
\coordinate (N1) at (1.6,1.6);

\draw[thick] (M) -- (M1) -- (M2) -- (M3) -- (N) -- (N1) -- cycle;

% labels
\node[above left] at (M) {$M$};
\node[above right] at (M1) {$M_1$};
\node[right] at (M2) {$M_2$};
\node[left] at (M3) {$M_3$};
\node[below] at (N) {$N$};
\node[left] at (N1) {$N_1$};

\end{tikzpicture}

（1）写出点 $ P $ 到居民区 $ A $ 的“$ L $ 路径”长度最小值的表达式（不要求证明）；

（2）若以原点 $ O $ 为圆心，半径为 $ 1 $ 的圆的内部是保护区，“$ L $ 路径”不能进入保护区，请确定点 $ P $ 的位置，使其到三个居民区的“$ L $ 路径”长度之和最小.
\end{problem}

\begin{answer}

\end{answer}

\begin{solutions}

\end{solutions}